\chapter{Il tessuto nervoso}

% ---------------------------------------------------------------------------
\section{Generalità}

I tessuti sono gruppi di cellule che lavorano in cooperazione ed hanno simili caratteristiche
morfologiche, permettendo una ``divisione del lavoro'' fra le cellule dell'organismo. Nel corpo
umano si distinguono 4 tipi di tessuto: epiteliale, connettivo, muscolare, nervoso.

\textbf{Tessuto nervoso}: due popolazioni cellulari $\rightarrow$ \textbf{neuroni} ($\sim$100
miliardi) e \textbf{cellule della glia} ($\sim$1000 miliardi). Non contiene matrice extracellulare.
\begin{itemize}
  \item \textbf{SNC} (sistema nervoso centrale): cervello e midollo spinale;
  \item \textbf{SNP} (sistema nervoso periferico): nervi e gangli.
\end{itemize}

% ---------------------------------------------------------------------------
\section{Il neurone}

\textbf{Neurone}: cellula eccitabile $\rightarrow$ risponde a stimoli fisici/chimici variando le
concentrazioni ioniche ai due lati della membrana. Il segnale = variazione del \textbf{potenziale
d'azione} (trasduzione), propagato lungo la membrana. Trasmette ad altre cellule tramite
\textbf{sinapsi}; dotato di \textbf{memoria}.
\begin{itemize}
  \item come ogni cellula: membrana, nucleo, mitocondri, sintesi proteica, produzione di energia;
  \item lo distinguono i prolungamenti: \textbf{dendriti} (portano info verso il corpo) e
    \textbf{assoni} (le portano lontano); contiene \textbf{neurotrasmettitori}.
\end{itemize}

\subsection{Le quattro regioni funzionali}
\begin{itemize}
  \item \textbf{dendriti}: raccolgono i segnali; sulla superficie le \textbf{spine dendritiche};
  \item \textbf{corpo cellulare}: nucleo (con nucleolo), mitocondri, ribosomi, organuli;
    \textbf{corpi di Nissl} e \textbf{neurofilamenti};
  \item \textbf{assone}: origina dal \textbf{monticolo assonico} (con \textbf{segmento iniziale}),
    conduce il potenziale d'azione alle terminazioni; può essere mielinizzato;
  \item \textbf{bottoni sinaptici}: trasmettono l'impulso ad altri neuroni o a effettori (muscoli,
    ghiandole).
\end{itemize}
Nel citoplasma: Golgi, REL e RER, poliribosomi e ribosomi liberi, microtubuli e microfilamenti
(anche nell'assone). Il bottone sinaptico contiene vescicole con neurotrasmettitore, rilasciato
nello spazio sinaptico e captato dai recettori postsinaptici. Lungo l'assone la guaina mielinica
(cellule di Schwann) è interrotta dai \textbf{nodi di Ranvier}; possibili sinapsi assone-assone.

% ---------------------------------------------------------------------------
\section{Classificazione strutturale}

Quattro tipi per morfologia:
\begin{itemize}
  \item \textbf{anassonico}: privo di assone identificabile, solo dendriti;
  \item \textbf{bipolare}: un dendrite e un assone;
  \item \textbf{pseudounipolare}: unico prolungamento che si divide a T (segmento iniziale +
    assone); detto anche bipolare o ``a T'';
  \item \textbf{multipolare}: vari dendriti e un unico assone.
\end{itemize}
Tutti (tranne gli anassonici) hanno un solo assone.

I neuroni possono anche essere classificati in base all'aspetto formale (cellule piramidali,
stellate, fusiformi) e in base al neurotrasmettitore utilizzato.

\subsection{Osservazioni istologiche}
\begin{itemize}
  \item midollo allungato: \textbf{neuroni multipolari} con zona tigroide del Nissl nel pirenoforo;
    nucleoli nei nuclei tagliati al centro; nuclei piccoli e rotondi = neuroglia (Nissl, 200x);
  \item corteccia telencefalica: neurone multipolare (colorazione di Cajal, 200x).
\end{itemize}

% ---------------------------------------------------------------------------
\section{Classificazione funzionale}

Organizzazione: recettori $\rightarrow$ SNP $\rightarrow$ SNC.
\begin{itemize}
  \item \textbf{recettori}: \textbf{interocettori}, \textbf{esterocettori}, \textbf{propriocettori}
    $\rightarrow$ inviano info via \textbf{fibre afferenti} ai \textbf{neuroni sensitivi} dei gangli
    periferici $\rightarrow$ al SNC;
  \item \textbf{interneuroni} (nel SNC): coordinano gli impulsi motori;
  \item sistema motorio:
    \begin{itemize}
      \item \textbf{somatico}: neuroni motori somatici $\rightarrow$ (fibre efferenti)
        $\rightarrow$ muscolatura scheletrica;
      \item \textbf{viscerale}: neuroni motori viscerali del SNC $\rightarrow$ (fibre pregangliari)
        $\rightarrow$ neuroni del ganglio periferico $\rightarrow$ (fibre postgangliari)
        $\rightarrow$ effettori viscerali (muscolatura liscia, ghiandole, muscolatura cardiaca,
        cellule adipose).
    \end{itemize}
\end{itemize}

% ---------------------------------------------------------------------------
\section{La sinapsi}

\textbf{Sinapsi}: giunzione che connette i neuroni e trasmette il segnale propagato lungo l'assone.
Per bersaglio:
\begin{itemize}
  \item \textbf{asso-somatica}: assone $\rightarrow$ corpo cellulare;
  \item \textbf{asso-dendritica}: assone $\rightarrow$ dendrite;
  \item \textbf{asso-assonica}: assone $\rightarrow$ assone.
\end{itemize}

\subsection{Sinapsi elettriche e chimiche}
\begin{itemize}
  \item \textbf{elettriche} (\textit{gap junctions}): i citoplasmi comunicano tramite canali di
    giunzione $\rightarrow$ passaggio di ioni; conduzione in \textbf{entrambe} le direzioni;
  \item \textbf{chimiche}: vescicole sinaptiche fondono con la membrana presinaptica $\rightarrow$
    neurotrasmettitore nello spazio sinaptico (\textit{synaptic cleft}) $\rightarrow$ recettori
    postsinaptici $\rightarrow$ canali ionici; conduzione in \textbf{una sola} direzione.
\end{itemize}
Nella sinapsi chimica lo spazio contiene proteine di regolazione della giunzione e del metabolismo
del neurotrasmettitore. Ogni vescicola $\sim$10.000 molecole di neurotrasmettitore.

\subsection{Sinapsi vescicolari}
Effetti postsinaptici brevi: il neurotrasmettitore è rimosso per via enzimatica o riassorbito. Sul
corpo del neurone: numerosi \textbf{bottoni sinaptici} connessi dall'\textbf{arborizzazione
terminale} agli assoni (rivestimento mielinico); visibili \textbf{dendriti} e \textbf{processi
gliali} tra le sinapsi.

% ---------------------------------------------------------------------------
\section{Circuiti neuronali}
\begin{itemize}
  \item \textbf{divergente}: un neurone stimola un numero crescente di neuroni a valle (stimolazione
    diffusa);
  \item \textbf{convergente}: più neuroni convogliano su un singolo neurone;
  \item \textbf{elaborazione seriale}: neuroni/gruppi in sequenza;
  \item \textbf{elaborazione in parallelo}: elaborazione simultanea su più vie;
  \item \textbf{riverberante}: feedback (eccitatorio o inibitorio), l'impulso torna via collaterale
    sul neurone di partenza.
\end{itemize}

% ---------------------------------------------------------------------------
\section{Le cellule della glia (nevroglia)}

\rubrica{Nel SNP}
\begin{itemize}
  \item \textbf{cellule satelliti}: circondano i corpi neuronali nei gangli; regolano $CO_2$,
    $O_2$, nutrienti e neurotrasmettitori;
  \item \textbf{cellule di Schwann}: circondano gli assoni del SNP, li mielinizzano, partecipano
    alla riparazione dopo lesione.
\end{itemize}
\rubrica{Nel SNC}
\begin{itemize}
  \item \textbf{oligodendrociti}: mielinizzano gli assoni del SNC + supporto strutturale;
  \item \textbf{astrociti}: barriera ematoencefalica, supporto, regolazione di ioni/nutrienti/gas,
    riciclo dei neurotrasmettitori, cicatrici dopo lesione;
  \item \textbf{microglia}: rimuove detriti, rifiuti e patogeni per fagocitosi;
  \item \textbf{cellule ependimali}: rivestono ventricoli e canale centrale; regolano produzione,
    circolazione e riassorbimento del liquido cerebrospinale (LCR).
\end{itemize}
Dettagli:
\begin{itemize}
  \item \textbf{ependimali}: epitelio monostratificato delle cavità del SNC (ventricoli, acquedotto
    del mesencefalo, canale midollare), con microvilli e ciglia unite da giunzioni;
  \item \textbf{microglia}: ``colla'' del sistema nervoso $\rightarrow$ regola comunicazione
    sinaptica e crescita neuronale, oltre ai processi immunitari;
  \item \textbf{astrociti}: le più numerose; forma stellata, espansioni laminari sui corpi
    neuronali, dendriti e sinapsi (sostanza grigia) e tra i fasci mielinizzati (sostanza bianca).
\end{itemize}

% ---------------------------------------------------------------------------
\section{La mielina}

\textbf{Mielina}: migliora la conduzione dell'impulso; struttura multilamellare biancastra. La
mielinizzazione inizia nel periodo embrionale e si completa in tempi diversi.
\begin{itemize}
  \item nelle fibre mielinizzate l'impulso è più veloce (conduzione ``saltatoria'', da nodo a
    nodo); le amieliniche sono di calibro inferiore;
  \item \textbf{oligodendrociti}: fino a 60 assoni; \textbf{cellule di Schwann}: un solo assone;
  \item guaina = composizione della membrana degli oligodendrociti/Schwann $\rightarrow$ colore
    bianco = \textbf{sostanza bianca};
  \item nel SNP: ogni cellula di Schwann fa un segmento internodale di un solo assone; al nodo, i
    lembi di due Schwann successive ricoprono l'assone isolandolo;
  \item composizione: 40\% acqua; della parte solida, 80\% lipidi (colesterolo, fosfolipidi) +
    proteine e carboidrati;
  \item struttura (esterno$\rightarrow$interno): fibre collagene dell'endonevrio, lamina basale,
    guaina mielinica, nucleo della cellula di Schwann, nodo di Ranvier, assone.
\end{itemize}

\subsection{Velocità di conduzione}
Con la mielinizzazione il segnale viaggia fino a 150 m/s; senza conduzione saltatoria si ridurrebbe
a $\sim$5 m/s.

\subsection{Il processo di mielinizzazione}
Un assone ``nudo'' affonda nel citoplasma di una cellula di Schwann/oligodendrocita $\rightarrow$ un
lembo ruota avvolgendosi a spirale $\rightarrow$ sovrapposizioni di membrana plasmatica (il
citoplasma viene rimosso a ogni spira) $\rightarrow$ molti doppi strati fosfolipidici isolanti.
\begin{itemize}
  \item SNP: una Schwann mielinizza un solo assone, oppure racchiude senza mielinizzare numerosi
    assoni amielinici;
  \item SNC: un oligodendrocita avvolge e mielinizza più assoni contemporaneamente.
\end{itemize}

% ---------------------------------------------------------------------------
\section{L'assone e la velocità di conduzione}

\textbf{Assone}: conduttore centrifugo rispetto al corpo cellulare. La velocità aumenta col
diametro (minore resistenza assiale $\rightarrow$ maggiore costante di spazio), come nelle fibre
amieliniche. Assoni primari afferenti per diametro e velocità:
\begin{table}[htbp]
  \centering \small \renewcommand{\arraystretch}{1.2}
  \begin{tabular}{|c|c|c|}
    \hline
    \textbf{Tipo} & \textbf{Diametro} & \textbf{Velocità} \\
    \hline
    A$\alpha$ & 13--20\,$\mu$m & 80--120 m/s \\
    A$\beta$  & 6--12\,$\mu$m  & 35--75 m/s \\
    A$\delta$ & 1--5\,$\mu$m   & 5--35 m/s \\
    C (amielinica) & 0,2--1,5\,$\mu$m & 0,5--2,0 m/s \\
    \hline
  \end{tabular}
\end{table}

\subsection{Riassumendo}
\begin{itemize}
  \item \textbf{amieliniche}: assoni in una nicchia delle cellule di Schwann, senza rivestimento;
  \item \textbf{mieliniche}: cellule mieliniche avvolte più volte $\rightarrow$ manicotti che
    schermano dalla corrente ionica; interruzioni ai \textbf{nodi di Ranvier} $\rightarrow$
    \textbf{conduzione saltatoria}.
\end{itemize}

% ---------------------------------------------------------------------------
\section{L'assone gigante di calamaro}

Organismo modello per l'eccitabilità: il \textbf{calamaro}. Assone gigante (fino a 1 mm di
diametro) $\rightarrow$ propulsione a getto d'acqua. Scoperto nel 1909, ha permesso di registrare
velocità fino a 210 m/s (tra le più elevate mai misurate).

% ---------------------------------------------------------------------------
\section{Cenni storici: elettricità e stimolazione del sistema nervoso}

Il tessuto nervoso è \textbf{eccitabile}; la stimolazione elettrica/magnetica del SNC è nota
dall'antichità: già nel 43--48 a.C.\ greci e romani sfruttavano i fenomeni bioelettrici di alcuni
pesci per cefalea ed epilessia.

\textbf{Pesci elettrici} (pesce gatto elettrico, torpedine, anguilla elettrica, pesce del Nilo,
razza elettrica): generano campi elettrici tramite organi di origine muscolare (difesa dai
predatori). Un lavoro del 1972 (\textit{Biochemical Journal}, 128, 833--846) documentò
localizzazione, innervazione e struttura degli organi elettrici della torpedine. Una rassegna di
storia della neurologia (Maggioni et al., \textit{Neurological Sciences}, 2005, 26, S431--S433)
ricostruì le terapie per la cefalea di Scribonio Largo. Riferimento iconografico: C.\ Bischoff,
\textit{Commentario de Usu Galvanismi} (1801). Documentate anche apparecchiature per l'elettroshock
e la terapia elettroconvulsivante.

Evoluzione delle tecniche di neurostimolazione (rassegna Hoy \& Fitzgerald, \textit{Nature Reviews
Neurology}, 2010):
\begin{itemize}
  \item \textbf{ECT} (anni '30): crisi convulsiva indotta da corrente elettrica via elettrodi sullo
    scalpo; meccanismo non noto (ipotesi: correzione di neurotrasmettitori inibitori, neurogenesi
    ippocampale);
  \item \textbf{tDCS} (anni '60): debole corrente allo scalpo (anodo/catodo) $\rightarrow$ altera
    l'eccitabilità spostando il potenziale di membrana (depol./iperpol.); effetto fino a $\sim$1 h;
  \item \textbf{TMS} (anni '90): campo magnetico focale intenso ($\sim$2 T) e variabile $\rightarrow$
    campo elettrico corticale $\rightarrow$ depolarizzazione/scarica; treni ripetuti modificano
    l'attività;
  \item \textbf{DBS} (2005): elettrodi impiantati in regioni cerebrali, stimolazione ad alta
    frequenza; collegati a un generatore (pacemaker) sotto la cute del torace;
  \item \textbf{EpCS} (2007): stimolazione corticale diretta con elettrodi sopra la dura madre;
  \item \textbf{100 Hz MST} (2008): crisi focale da TMS ripetitiva ad alta frequenza; a differenza
    dell'ECT nessuna corrente attraversa le regioni profonde; meccanismi non noti.
\end{itemize}

% ---------------------------------------------------------------------------
\section{Applicazioni sportive: la tDCS nel salto con gli sci}

Caso applicativo: \textbf{tDCS} negli specialisti del salto con gli sci. Studio su \textit{Nature}
(17 marzo 2016, 531, 283), titolo giornalistico \textit{``Performance boost paves way for `brain
doping'''} $\rightarrow$ la stimolazione elettrica sembra migliorare la resistenza (studi
preliminari). Stimolazione della corteccia motoria (4 soggetti) vs placebo \textit{sham} (3
soggetti):
\begin{itemize}
  \item migliora le prestazioni;
  \item riduce la percezione della fatica.
\end{itemize}
Nel dettaglio: +70\% forza di stacco (\textit{jumping force}), +80\% coordinazione rispetto al
placebo (dispositivo ``Halo'', annunciato nel febbraio dello stesso anno).
